%% ma: L11-B17-0002
%% lop: 11
%% chuong: 5
%% bai: 17
%% dang: 11.17.1
%% loai: TL
%% muc: TH
%% dap_an: "-"
%% nguon: "tai-lieu-hoc-tap-mon-toan-lop-11-theo-chuong-trinh-giao-duc-pho-thong-2018.pdf (Lớp Toán Thầy Duy, Phạm Lê Duy, THPT Chu Văn An), Chương 3 Bài 3, Dạng 1, Ví dụ 2 (tr.387-389)"
%% kien_thuc: "Lớp 11 Bài 17 (hàm số liên tục)"
%% trang_thai: chinh-thuc
%% ghi_chu: "Đã duyệt 03/11/2026."
\begin{ex}
Cho hàm số $f(x)=\begin{cases}x^2+1&\text{khi }x>0\\x&\text{khi }x\le0\end{cases}$. Xét tính liên tục của hàm số tại điểm $x_0=0$.
\loigiai{\sloppy $f(0)=0$; $\lim\limits_{x\to0^+}f(x)=\lim(x^2+1)=1$, $\lim\limits_{x\to0^-}f(x)=\lim x=0$. Hai giới hạn một bên khác nhau nên không tồn tại $\lim\limits_{x\to0}f(x)$; hàm số gián đoạn tại $x_0=0$.}
\end{ex}
